\section{A positive delayed-use experiment with exact causal risk}
\label{sec:v32-delayed}

The compatibility equations have an elementary test in which two schedules
are both admissible and local prediction costs are strictly positive.
The calculation below also records the effect of \emph{not} supplying a
persistent public seed. It is a separate finite positive experiment,
not a substitute for the general collision theorem.

Fix $0<\gamma\le1/8$ and $0<e\le \gamma/2$. Let
$Z_1,Z_2,V_1,V_2$ be independent symmetric signs. The four acquisition
reports $X_1,X_2,U_1,U_2$ are conditionally independent, with
\[
 \Pr(X_j=x\mid Z_j=z)=\tfrac12(1+xz/2),\qquad
 \Pr(U_j=u\mid V_j=v)=\tfrac12(1+uv/2).
\]
All acquisition probabilities are at least $1/4$, and the reports are
independent symmetric signs unconditionally. Query $j$ is a Bernoulli
trial of latent success probability
\[
             \tfrac12+2\gamma Z_j+2eV_j.
\]
Its success and failure probabilities are at least $1/8$. Once its two
reports have been acquired, the full-history conditional mean is
\begin{equation}\label{eq:v32-delayed-target}
              \tfrac12+\gamma X_j+eU_j.
\end{equation}

There are two public, deterministic admissible schedules. The
\emph{early} schedule reads all four acquisition reports as one block
before checkpoint 1; no informative acquisition occurs between the
checkpoints. The \emph{just-in-time} schedule reads $(X_1,U_1)$ before
checkpoint 1 and $(X_2,U_2)$ before checkpoint 2. In both schedules
memory is charged after each acquisition block and at each checkpoint;
there are two retained labels, the clock is free, no public seed is
retained, and no score or past prediction is available for later updates.
Each report is acquired exactly once. Only one independently selected
checkpoint query is executed. All four scheduled acquisition trials are
completed even when checkpoint 1 is scored, and the score is not fed back.
Thus both schedules use four acquisition reports and one scored query.
Queries at unselected checkpoints are not additional observations. The block-boundary charging
rule is part of this example; it is not the per-report rule of the scalar
monomial theorem.

Let $V_{\rm early}(e)$ and $V_{\rm jit}(e)$ denote the infimum of the
maximum of the two expected squared excesses over the full-history
predictors. Private update randomization is allowed in both infima.

\begin{theorem}[Exact unseeded delayed-use law]\label{thm:v32-delayed}
For the experiment and resource convention just specified,
\begin{equation}\label{eq:v32-delayed-exact}
 \begin{split}
 V_{\rm jit}(e)&=e^2,\\
 V_{\rm early}(e)&=\gamma^2+e^2-
       \max\left\{\frac{\gamma^2}{3},
                       \frac{(2\gamma+e)^2}{15}\right\}.
 \end{split}
\end{equation}
Both values are attained by deterministic two-label encoders and fixed
stage-dependent decoders. In particular,
\[
 \frac{V_{\rm early}(e)}{V_{\rm jit}(e)}\longrightarrow\infty
                         \quad\hbox{as }e\downarrow0.
\]
This comparison is between two admissible common-controller designs with
the same total acquisition count, the same query protocol and the same
retained-label and seed convention. It is not a comparison with
incompatible checkpoint-specific initial designs.
\end{theorem}

\begin{proof}
We work with the centered targets
$Y_j=\gamma X_j+eU_j$, whose variances are $\gamma^2+e^2$.
For a single current target the four values are
$-\gamma-e,-\gamma+e,\gamma-e,\gamma+e$, all with probability $1/4$.
For squared loss an optimum two-center scalar quantizer partitions an
ordered set into contiguous cells: assigning each point to a nearest
center cannot increase loss. The middle split has distortion $e^2$.
The two extreme splits have explained variance $(\gamma+e)^2/3$, whereas
the middle split has explained variance $\gamma^2$. Since
$e\le \gamma/2$ implies $(\gamma+e)^2/3\le3\gamma^2/4$, the middle split is
optimal. A fresh current pair is independent of all earlier retained
information, which cannot enlarge the two available reconstruction
values at that checkpoint. Resetting to the sign $X_j$ therefore attains
$e^2$ at both checkpoints, and the one-block lower bound proves the
just-in-time formula.

For the early schedule, later private transitions without new information
cannot improve on decoding directly from the first retained label.
Indeed, conditional on that label, averaging any later random prediction
cannot increase squared loss, and its conditional expectation is a
function of that label alone. Thus every early controller is dominated
by a stochastic binary encoder of the four-report word followed by two
fixed decoders.

Write $Q\in[0,1]$ for the conditional probability of label 1 given that
word, let $p=\mathbb E Q$, and put $v_j=\mathbb E(Y_jQ)$. If
$0<p<1$, centroid decoding makes checkpoint $j$'s risk
\begin{equation}\label{eq:v32-binary-centroid}
 \gamma^2+e^2-\frac{v_j^2}{p(1-p)}.
\end{equation}
A constant encoder explains no variance. Independent reflection of either
coordinate pair $(X_j,U_j)$ changes the sign of $v_j$ without changing
$p$ or the objective. We may therefore take $v_1,v_2\ge0$.
Replacing $Q$ by the average of itself and its coordinate-swapped version
preserves $p$ and changes both $v_j$ to $(v_1+v_2)/2$. This cannot
decrease their smaller square. It is an allowed private randomized
encoder, not a public mixture of decoders. We may consequently optimize
over encoders with $v_1=v_2$.

Put $S=Y_1+Y_2$. At a fixed $p$, the largest possible value of
$\mathbb E(SQ)$ is obtained by taking an upper tail of $S$, randomizing
only at a tied threshold. This follows by transferring any positive
assignment mass from a smaller value of $S$ to a larger one. Symmetric
assignment within ties gives equal coordinate correlations. Relabeling
and joint reflection reduce the range to $0<p\le1/2$.
On an interval between successive tail masses, write
$\mathbb E(SQ)=A+sp$, where $s\ge0$ and $A\ge0$. Apart from a
positive constant, the explained variance is
\[
                       \frac{(A+sp)^2}{p(1-p)}.
\]
Its derivative has the sign of $(s+2A)p-A$ whenever its numerator is
nonzero. It therefore has no interior strict maximum. It is enough to
consider the finite tail endpoints, including a half-mass endpoint.

The sixteen values of $S$ are sums of
$\gamma(X_1+X_2)+e(U_1+U_2)$. For $e\le \gamma/2$, the sums of the
largest $n$ values, $1\le n\le8$, are the following $A_n$:
\begin{center}
\begin{tabular}{c c c}
\toprule
$n$ & $A_n$ & explained variance $A_n^2/[4n(16-n)]$\\
\midrule
1 & $2\gamma+2e$ & $(\gamma+e)^2/15$\\
2 & $4\gamma+2e$ & $(2\gamma+e)^2/28$\\
3 & $6\gamma+2e$ & $(3\gamma+e)^2/39$\\
4 & $8\gamma$ & $\gamma^2/3$\\
5 & $8\gamma+2e$ & $(4\gamma+e)^2/55$\\
6 & $8\gamma+4e$ & $(2\gamma+e)^2/15$\\
7 & $8\gamma+4e$ & $4(2\gamma+e)^2/63$\\
8 & $8\gamma+4e$ & $(2\gamma+e)^2/16$\\
\bottomrule
\end{tabular}
\end{center}
Ties at $e=\gamma/2$ do not change these sums. Rows 1 and 2 are bounded
by row 6, and rows 7 and 8 are also bounded by row 6. Row 3 is bounded
by row 4, since
$4\gamma^2-6\gamma e-e^2\ge3\gamma^2/4$.
For row 5, if $e\le \gamma/4$ then
$7\gamma^2-24\gamma e-3e^2\ge13\gamma^2/16>0$, so row 4 dominates.
If $e\ge \gamma/4$, then
$-\gamma^2+5\gamma e+2e^2\ge3\gamma^2/8>0$, so row 6 dominates.
Thus the largest explained variance is the maximum in
\eqref{eq:v32-delayed-exact}.

Both remaining candidates are attained without encoder randomization.
The four-word tail is $\{X_1=X_2=1\}$. Its two conditional centered
means are $(\gamma,\gamma)$ and $(-\gamma/3,-\gamma/3)$.
The six-word tail adds
$\{X_1=-X_2,\ U_1=U_2=1\}$. Its centered means in either coordinate
are $(2\gamma+e)/3$ and $-(2\gamma+e)/5$. Add $1/2$ to obtain the allowed
prediction probabilities. These two encoders attain the computed maximum
of the explained variances, proving the exact formula. Finally,
$V_{\rm early}(e)\to2\gamma^2/3>0$ while $V_{\rm jit}(e)=e^2$.
\end{proof}

The active branch changes at $e=(\sqrt5-2)\gamma$. The formula is an
analytic solution of the private-randomization minimax problem, not merely
an enumeration of deterministic partitions. A retained public coin gives
a different decision problem and need not have this value. For example,
a fair public coin that selects which large sign $X_j$ to retain gives
risk $\gamma^2/2+e^2$ at each checkpoint. This observation is an admissible
upper bound only in that enlarged resource convention. It is not used in
Theorem~\ref{thm:v32-delayed}.

At each checkpoint the currently scored marginal law is the same in the
two schedules, and its two-label distortion is $e^2$. Nevertheless, the
common early controller must preserve competing directions for later
use. Equations~\eqref{eq:v32-compatible} retain that constraint. The
example exhibits its consequence in a model with positive likelihoods and
positive local risks; the uniform approximation and relative certificate
results address the whole prescribed monomial family, rather than this
one finite test alone.
